% ============================================================================
% localization-guide-EN.tex — SUIT Localization Guide (English)
% ============================================================================
%
% STATUS: populated. One row per parameterize entry.
%
% ----------------------------------------------------------------------------
% PURPOSE
% ----------------------------------------------------------------------------
% The SUIT documents are written to be adopted directly by any university. The
% body of every document uses role-based, vendor-neutral, institution-neutral,
% and jurisdiction-neutral language (see .claude/rules/suit-domain.md). Wherever
% a concrete local anchor would otherwise be hard-coded into the prose -- a
% national CSIRT, a data-protection authority, a national NREN, a specific
% governance body, a national transposition of an EU instrument -- the body
% already reads in a GENERIC ROLE TERM. This guide tells an adopting institution
% which concrete local value to read each generic role term as.
%
% ----------------------------------------------------------------------------
% THE HYBRID LOCALIZATION MODEL  (read this first)
% ----------------------------------------------------------------------------
% SUIT uses a two-tier ("hybrid") model rather than a single substitution sweep:
%
%   1. GENERIC BODY (default, always shipped)
%      The published SUIT documents read entirely in generic terms. They are
%      complete and correct as-is: an adopter can read and reason about them
%      without supplying any local value. The generic term is the canonical
%      reference; localization is additive, never required for comprehension.
%
%   2. LOCAL INSTANTIATION (per-adopter reading)
%      An adopting institution maps each generic role below to its own concrete
%      value (its NREN, its national CSIRT, its DPA, its supervisory authority
%      under the national NIS2 transposition, its highest governance body, etc.)
%      and reads the generic prose against that mapping, producing an
%      institution-specific understanding WITHOUT forking the generic source.
%
% IMPORTANT -- THIS IS DOCUMENTATION, NOT A BUILD-TIME OVERLAY.
%   * The SUIT document bodies contain FINISHED GENERIC PROSE. There are NO
%     mustache tokens, no {{PLACEHOLDERS}}, and no \locvar/\locdefine macros in
%     the source to expand. Nothing is substituted at compile time.
%   * The table below is a HUMAN-FACING reference: it tells an adopter, in plain
%     reading terms, "where the documents say X (generic role), read it as your
%     Y (local value), and here is where X appears." Localizing is an editorial
%     act performed by the adopter against their own jurisdiction and estate, not
%     an automated token sweep.
%   * The two tiers share ONE source. The generic edition is the source of truth;
%     a localized edition is the generic edition read (and, if an adopter
%     chooses, edited) against this lookup table.
%
% ----------------------------------------------------------------------------
% HOW TO READ THE TABLE
% ----------------------------------------------------------------------------
% Each row corresponds to one localization variable. The columns are:
%
%   VARIABLE        a stable handle for the localization point (for adopter
%                   notes / change tracking; it is NOT a token in the prose).
%   GENERIC ROLE    the generic role phrase as it appears in the SUIT body.
%   ADOPTER VALUE   what an adopting university reads it as / substitutes it with.
%   WHERE / NOTES   which document(s) it appears in and any scoping guidance,
%                   including whether the generic phrasing should simply be kept.
%
% Document keys used in the WHERE column:
%   SUITPOLICY         -- suit-policy/policy-EN.tex                (the Policy)
%   SUITPOLICYSUMMARY  -- suit-policy/companion/policy-summary-EN.tex (Summary)
%   SUITSOLUTION       -- suit-solution/solution-EN.tex            (Technical Solution)
%   CONSOLIDATION      -- ul-itsec-consolidation-EN.tex (out of suite scope; noted)
%
% "Keep" in the ADOPTER VALUE column means the generic phrasing is already
% institution-/jurisdiction-agnostic and an adopter normally leaves it unchanged;
% the WHERE / NOTES column states the condition under which it would still change.
%
% ============================================================================

\begin{center}
\small
\begin{longtable}{@{}p{0.20\linewidth}p{0.24\linewidth}p{0.24\linewidth}p{0.28\linewidth}@{}}
\caption{SUIT localization map — read each generic role as your local value.}
\label{tab:suit-localization}\\
\toprule
\textbf{Variable} & \textbf{Generic role (as it appears)} & \textbf{Adopter substitutes with} & \textbf{Where it appears / notes}\\
\midrule
\endfirsthead
\multicolumn{4}{@{}l}{\small\itshape Table~\ref{tab:suit-localization} (continued)}\\
\toprule
\textbf{Variable} & \textbf{Generic role (as it appears)} & \textbf{Adopter substitutes with} & \textbf{Where it appears / notes}\\
\midrule
\endhead
\midrule
\multicolumn{4}{r@{}}{\small\itshape continued on next page}\\
\endfoot
\bottomrule
\endlastfoot

% --- amount ------------------------------------------------------------------
\texttt{NIS2\_SANCTION\_CEILING} &
The generic NIS2 essential-entity sanction ceiling (EUR amount) &
Keep --- this is the generic NIS2 essential-entity ceiling, not a Luxembourg-specific figure &
SUITPOLICY (3 occurrences). Generic regulatory anchor; left unchanged unless the adopter's regime sets a different ceiling.\\
\midrule

% --- artifact ----------------------------------------------------------------
\texttt{INSTITUTION\_NAME} (bib) &
UL-suite bibliography path and citation keys &
The adopter's own bibliography; flag UL-specific keys (e.g.\ \texttt{unilu-amenagements}, \texttt{ilr-ul-essential}) for removal or parameterization &
SUITPOLICY, SUITPOLICYSUMMARY (5). These keys support home-institution facts and have no generic equivalent; remove/parameterize when localizing to the adopter's institution.\\
\midrule

% --- cost / staffing ---------------------------------------------------------
\texttt{STAFFING\_PROFILE} &
Staffing / FTE envelopes, cadences and windows &
The adopter's own staffing envelope; drop the wealth-of-country claim and the "excellence as sole objective" framing (or generalize to "an institution that sets excellence as a strategic objective") &
SUITPOLICY, SUITSOLUTION (24). Largest localization point; figures are illustrative of the home institution.\\
\midrule

% --- facilities --------------------------------------------------------------
\texttt{NATIONAL\_HPC\_FACILITY} &
National/regional supercomputer (MeluXina/EuroHPC) &
"A national/EuroHPC supercomputer", or the adopter's regional EuroHPC system &
SUITPOLICY, SUITSOLUTION (8).\\
\midrule

\texttt{IDENTITY\_FEDERATION\_PROFILE} &
Identity interfederation (eduGAIN) &
Keep for an EU edition; for a jurisdiction-agnostic edition gate the EU-law layer behind a "regulatory regime" parameter (EU as default). Parameterize the EU-first provider examples (Mistral/OVH/Scaleway), which are region-bound &
SUITPOLICY, SUITSOLUTION (6).\\
\midrule

\texttt{INSTITUTION\_HPC\_FACILITY} &
Institutional HPC facility/team (ULHPC) &
"The institution's HPC facility / HPC team" &
SUITPOLICY (21).\\
\midrule

\texttt{NREN\_NAME} &
National research \& education network (RESTENA) &
"The institution" / "our institution" --- this whole Q\&A is a Luxembourg-specific local-deployment objection &
SUITPOLICY, SUITPOLICYSUMMARY, SUITSOLUTION (41). Highest-frequency point.\\
\midrule

\texttt{NREN\_CSIRT} &
NREN security team (RESTENA-CSIRT) &
"The national NREN" (National Research and Education Network) and "the NREN CSIRT" &
SUITPOLICY, SUITSOLUTION (14).\\
\midrule

\texttt{RESEARCH\_BACKBONE\_SCALE} &
Region-bound research-backbone scale and capacity figures (G\'EANT/RENATER) &
Keep the campus-connection capacity range (10--100\,Gbps) as a jurisdiction-neutral engineering figure; parameterize the G\'EANT-specific membership and backbone-capacity numbers to "the regional research backbone (membership and backbone capacity per region)" &
SUITPOLICY (1). Affects 40+ NRENs, 50M users, 10,000+ institutions, 400\,Gbps wavelengths, multi-100\,Gbps national backbones.\\
\midrule

\texttt{INTERNAL\_EMERGENCY\_INCIDENT\_CONTACT} &
Institution-specific internal emergency incident-reporting number (9911) &
A generic "internal emergency / incident-reporting contact"; keep the operational pattern (mandatory reporting to ISO/CISO with an internal emergency line) &
CONSOLIDATION (1). Out of suite file scope --- no work-item; listed for completeness.\\
\midrule

% --- governance bodies -------------------------------------------------------
\texttt{ACADEMIC\_GOVERNING\_BODIES} &
Academic deliberative bodies (Faculty/University Council) &
Role-based collegial-body terms for academic deliberative councils &
SUITPOLICY (2).\\
\midrule

\texttt{EXECUTIVE\_LEADERSHIP} &
Executive leadership (Rectorate/Rector) &
A role-based term for the institution's central executive leadership &
SUITPOLICY, SUITSOLUTION (4).\\
\midrule

\texttt{GOVERNING\_BOARD} &
Top governing board (Board of Governors) &
Keep --- already largely role-based ("governing bodies", "board", "CIO"); confirm these remain generic governance roles rather than a specific board-of-governors/rectorate shape &
SUITPOLICY, SUITPOLICYSUMMARY, SUITSOLUTION (19).\\
\midrule

% --- pilot entities ----------------------------------------------------------
\texttt{PILOT\_ENTITIES} &
Pilot entity set &
Replace ULHPC with the role "the institution's HPC facility / team" &
SUITPOLICY, SUITPOLICYSUMMARY, SUITSOLUTION (6).\\
\midrule

% --- institution identity / profile -----------------------------------------
\texttt{INSTITUTION\_NAME} &
Subject institution name (University of Luxembourg) &
"The institution" / "a small institution" as the subject of the comparison objection &
SUITPOLICY, SUITSOLUTION (22).\\
\midrule

\texttt{INSTITUTION\_TYPE} &
Reference institution type/profile (R1 / research-intensive) &
Keep --- "research-intensive institution" is hedged with "typical"; note as a scoping assumption a non-research (teaching-led or administrative) profile would adjust &
SUITSOLUTION (9).\\
\midrule

\texttt{CS\_RESEARCH\_UNIT} &
Computer Science department (pilot/host unit) &
Role-based "entities with the strongest domain expertise (e.g., a computer-science department, a security/research centre, and the HPC team)" &
SUITPOLICY, SUITSOLUTION (23).\\
\midrule

\texttt{SECURITY\_RESEARCH\_CENTRE} &
Security research centre (SnT) &
"A cybersecurity-focused interdisciplinary research centre" &
SUITPOLICY (5).\\
\midrule

% --- macro -------------------------------------------------------------------
\texttt{INSTITUTION\_MACRO\_SET} &
\texttt{\textbackslash institution\{\}} macro set expanding to UL terms &
Keep --- already framed conditionally ("for institutions whose \dots\ framework is built on ITIL or COBIT"); institution-agnostic by construction &
SUITSOLUTION (1).\\
\midrule

% --- legal -------------------------------------------------------------------
\texttt{AI\_REGULATION} &
AI-governance regime (EU AI Act) &
For a jurisdiction-agnostic template, parameterize the legal-framework layer to the adopter's jurisdiction; the whole legal scaffolding (GDPR, NIS2, AI Act, EU Charter, CJEU/ECtHR) presumes an EU/European jurisdiction &
SUITPOLICYSUMMARY, SUITSOLUTION (8).\\
\midrule

\texttt{CER\_NATIONAL\_TRANSPOSITION} &
National CER transposition law (PL 8307) &
"The national CER (Critical Entities Resilience) transposition law"; keep CER Directive (EU) 2022/2557 as the generic anchor &
SUITPOLICY (6).\\
\midrule

\texttt{DATA\_PROTECTION\_REGIME} &
Data-protection regime (GDPR) &
Keep US-funding references where they describe external precedent institutions; where 800-171/CUI is presented as a driver, frame as "where the institution processes regulated data (e.g.\ controlled, health, classified)" &
SUITPOLICY, SUITSOLUTION (7).\\
\midrule

\texttt{EID\_TRUST\_SERVICES\_REGIME} &
eID / trust-services regime (eIDAS\,2) &
Frame the regulatory anchors as "the institution's applicable data-protection, network-and-information-security, and trust-services regime", giving GDPR/NIS2/eIDAS\,2 as the worked instantiation; keep ISO\,27001 and NIST CSF (jurisdiction-neutral) &
SUITSOLUTION (4).\\
\midrule

\texttt{NATIONAL\_EQUAL\_TREATMENT\_LAW} &
National equal-treatment / anti-discrimination law (Law of 28 Nov 2006) &
A role-based reference to "the national transposition of Directive 2000/78/EC into labour law" &
SUITPOLICY (8).\\
\midrule

\texttt{ESSENTIAL\_ENTITY\_CLASSIFICATION\_BASIS} &
Essential-entity classification status \& basis &
Generalize "the University [is] an essential entity" to "institutions under the NIS2 essential-entity (or important-entity) regime", making classification a configurable premise &
SUITPOLICY, SUITSOLUTION (11).\\
\midrule

\texttt{NATIONAL\_CYBER\_STRATEGY} &
National cybersecurity strategy (NCSS III) &
A role-based reference to the national cybersecurity strategy's academic-cooperation objective &
SUITPOLICY (4).\\
\midrule

\texttt{CYBERSECURITY\_BASELINE\_REGIME} &
Cybersecurity baseline regime (NIS2) &
Generalize "including Luxembourg" to "including the institution's jurisdiction"; frame the EU/CoE legal order as "the applicable supranational and national legal framework" so non-EU adopters can substitute their own &
SUITPOLICY, SUITSOLUTION (15).\\
\midrule

\texttt{NIS2\_NATIONAL\_TRANSPOSITION} &
National NIS2 transposition law (PL 8364) &
Generalize the lawful-basis argument to "a public-sector or public-interest institution"; note the NIS2 statutory duty depends on national transposition (stated without naming the transposing law/country) &
SUITPOLICY, SUITPOLICYSUMMARY, SUITSOLUTION (15).\\
\midrule

\texttt{REGULATORY\_REGIME} &
Composite regulatory regime / national transposition &
"Binding law in the EU and Council of Europe member states (including the institution's jurisdiction)" &
SUITPOLICY, SUITSOLUTION (10).\\
\midrule

\texttt{NATIONAL\_PUBLIC\_SECTOR\_ISP} &
State public-sector information-security policy (PSI-LU) &
A role-based reference to the national public-sector information-security baseline policy &
SUITPOLICY (4).\\
\midrule

\texttt{CROSS\_BORDER\_TRANSFER\_REGIME} &
Cross-border data-transfer regime (Schrems II) &
Keep --- the list is supranational (EU/international), framed generically for "a research-intensive university", with no national transposition named; already institution-agnostic &
SUITPOLICY, SUITSOLUTION (11).\\
\midrule

\texttt{NATIONAL\_CYBER\_COMMISSARIAT\_FOUNDING\_LAW} &
National founding statute of the cyber-resilience commissariat (loi du 23 juillet 2016) &
Parameterize to the national founding statute of the cyber-resilience/protection commissariat; keep the role (agency exists by national law, sits under the executive), drop the 23 July 2016 date and the loi-modifi\'ee form &
SUITPOLICY, SUITSOLUTION (2).\\
\midrule

% --- jurisdiction ------------------------------------------------------------
\texttt{DATA\_CUSTODY\_SOVEREIGNTY\_PROFILE} &
On-premises key custody / no-third-country-disclosure as the structural data-sovereignty optimum &
Generalize the data-dimension optimum to "custody arrangement that minimises the applicable disclosure-regime exposure" rather than hard-coding on-premises custody; gate the "no third-country disclosure regime" criterion behind the adopter's jurisdiction profile &
SUITSOLUTION (3).\\
\midrule

\texttt{NATIONAL\_LEGISLATURE} &
National legislative chamber as deposit / enactment venue (Chambre des D\'eput\'es) &
A role-based national-legislature reference; keep that the instrument is national law passed by the national parliament, drop the chamber proper noun, the Grand-Duch\'e framing and the vote-tally specifics &
SUITPOLICY, SUITSOLUTION (2).\\
\midrule

% --- method ------------------------------------------------------------------
\texttt{NATIONAL\_RISK\_METHOD} &
National risk-analysis method (MONARC) &
Keep MONARC only as one of several examples, or parameterize the leading example by jurisdiction (listing it first implies a Luxembourg/ANSSI-LU context) &
SUITSOLUTION (1).\\
\midrule

% --- person ------------------------------------------------------------------
\texttt{AUTHOR\_NAME} &
Named author / contact person &
Keep --- these are public case-law attributions, not institution-specific persons; they generalize to any EU adopter &
SUITPOLICY, SUITSOLUTION (3).\\
\midrule

\texttt{AUTHOR\_AFFILIATION} &
Author-to-affiliation binding (named contact + UL host unit + uni.lu email) &
Parameterize the binding everywhere it appears: self-citation \texttt{@techreport} institution/note fields (\texttt{ng-tech-ref-2026}, \texttt{ng-dossier-2026}), the title-page contact-point block, the book front-matter "Citation:" line, and the PDF \texttt{\textbackslash hypersetup} pdfauthor/pdftitle metadata. The uni.lu email domain must be parameterized, not carried into another institution's edition; remove the pdfauthor AI-tool string &
SUITPOLICY, SUITSOLUTION (4).\\
\midrule

% --- regulators --------------------------------------------------------------
\texttt{NATIONAL\_CYBER\_COMPETENCE\_CENTRE} &
National cyber competence centre (NC3 / House of Cybersecurity) &
"An ISO/IEC 27005-aligned risk-analysis method/tool (e.g.\ a nationally provided open-source methodology)"; keep ISO/IEC 27005 as the generic anchor &
SUITPOLICY (5).\\
\midrule

\texttt{NATIONAL\_DPA} &
National data-protection authority (CNPD) &
"The competent data protection authority (DPA)" &
SUITPOLICY, SUITSOLUTION (11).\\
\midrule

\texttt{NATIONAL\_GOV\_CSIRT} &
Governmental CSIRT (GOVCERT.LU) &
A role-based term for the national/governmental CSIRT applicable to the public sector &
SUITPOLICY (7).\\
\midrule

\texttt{NATIONAL\_INFOSEC\_AGENCY} &
National information-security agency (ANSSI-LU) &
A role-based term for the national information-security agency authoring the state IS policy &
SUITPOLICY, SUITSOLUTION (13).\\
\midrule

\texttt{NATIONAL\_NOTIFICATION\_PLATFORM} &
National incident-notification platform (SERIMA) &
A role-based term for the national NIS2 incident-notification/risk-analysis platform; drop the launch date &
SUITPOLICY, SUITPOLICYSUMMARY (8).\\
\midrule

\texttt{NATIONAL\_SPOC\_AUTHORITY} &
NIS2 national single point of contact (HCPN) &
Generalize to "the national cybersecurity strategy, where it identifies cooperation with the academic and research sphere as a priority" &
SUITPOLICY (4).\\
\midrule

\texttt{NIS2\_SUPERVISORY\_AUTHORITY} &
NIS2 supervisory authority (ILR) &
A role-based term for the national NIS2 supervisory authority &
SUITPOLICY, SUITPOLICYSUMMARY, SUITSOLUTION (25).\\
\midrule

\texttt{NATIONAL\_EXECUTIVE\_DEPARTMENT} &
National executive department hosting the cyber-resilience commissariat (Minist\`ere d'\'Etat) &
A role-based reference (the national executive department responsible for the protection commissariat); keep that the commissariat sits under the head-of-government's authority, drop the Minist\`ere d'\'Etat proper noun and Francophone form &
SUITPOLICY, SUITSOLUTION (2).\\
\midrule

% --- vendor ------------------------------------------------------------------
\texttt{INCUMBENT\_IDENTITY\_SUITE\_VENDOR} &
Incumbent proprietary identity suite (Entra ID / Conditional Access) &
Keep as a generally-recognised market example of a "pure-commercial identity provider" (illustrative, not prescriptive); note the "designated EU region" carries the EU-jurisdiction assumption &
SUITSOLUTION (8). Named products: Microsoft Entra ID, Conditional Access, Customer Lockbox, Double Key Encryption.\\

\end{longtable}
\end{center}

% ============================================================================
% STEP 0 — SELECT YOUR REGULATORY REGIME
% ============================================================================
%
% This prose step precedes the per-variable substitution above for any adopter
% whose binding regime is not the EU/EEA baseline. It points at the shared
% Regulatory Applicability module (shared/regulatory-applicability-EN.tex),
% which carries the size- and jurisdiction-invariant obligation spine (F1--F6),
% the obligation->mechanism matrix, the regime typology, and the derivation
% method for unlisted regimes. Selecting the regime here fixes the value of the
% legal/regulator localization variables in the table above (DATA_PROTECTION_
% REGIME, CYBERSECURITY_BASELINE_REGIME, AI_REGULATION, EID_TRUST_SERVICES_
% REGIME, CROSS_BORDER_TRANSFER_REGIME, NATIONAL_DPA, NIS2_SUPERVISORY_
% AUTHORITY, NATIONAL_GOV_CSIRT, and the *_NATIONAL_TRANSPOSITION rows) in one
% coherent move rather than row by row.

\bigskip
\noindent\textbf{Step~0 --- Select your regulatory regime (do this first).}
The legal and regulator rows of the table above (data-protection regime,
cybersecurity-baseline regime, AI-governance regime, eID/trust-services regime,
cross-border-transfer regime, and the national-authority and
national-transposition rows) are all instances of a single decision: which
regulatory regime binds the adopting institution. Make that decision once,
using the shared \emph{Regulatory Applicability} module
(\texttt{shared/regulatory-applicability-EN.tex}), and the individual rows
follow coherently. The module separates the part that does not change from the
part that does: the size- and jurisdiction-invariant functional-obligation
spine (F1--F6) and the obligation~$\rightarrow$~mechanism matrix are fixed for
every adopter; only the \emph{instrument} column --- the local statute,
competent authority and deadline --- is substituted per regime.

\smallskip
\noindent\emph{If your regime is one of the typology profiles}, read it
directly from the module's regime typology and adopt its F1--F6 instrument
column as your localization. The profiles shipped are:

\begin{itemize}[leftmargin=1.5em,nosep]
  \item \textbf{European Union / EEA} --- deep regime, the SUIT baseline anchor
        (GDPR, NIS2, EU AI Act, eIDAS\,2, GDPR Chapter~V transfers);
  \item \textbf{United States} --- deep regime, sectoral and fragmented (FERPA,
        HIPAA, GLBA/FTC Safeguards, NIST SP~800-171 / CMMC, EAR/ITAR,
        CCPA/CPRA);
  \item \textbf{United Kingdom} --- light regime, post-Brexit diverging (UK
        GDPR + DPA~2018 + DUAA~2025, NIS Regulations~2018, UK eIDAS, PSBAR);
  \item \textbf{Brazil} --- light regime, distributed federal stack (LGPD,
        E-Ciber, CTIR.Gov, the Inclusion Law, the Access to Information Law);
  \item \textbf{South Africa} --- light regime, single national jurisdiction
        (POPIA, the Cybercrimes Act, the ECT Act, PEPUDA / Equality Act).
\end{itemize}

\smallskip
\noindent\emph{If your regime is not listed}, derive your own regime row with
the module's seven-step mapping method
(\texttt{regulatory-applicability-EN.tex}, ``Mapping method: deriving a row for
an unlisted regime''): (1) identify your data-protection authority and its
security-of-processing duty (F1); (2) identify your lawful-basis and
risk-governance regime (F2); (3) record your shortest breach/incident
notification deadline (F3); (4) determine your essential-/important-entity
status and its basis (F2/F3 scope); (5) establish your cross-border transfer
regime (F4); (6) add your sector overlays --- trust/eID, AI governance, export
control, health, financial-aid, accessibility (F5/F6); (7) write each local
instrument, authority and deadline into the instrument column against the
matching F~row, leaving the SUIT control-mechanism column unchanged. The result
is your own regime profile in the same shape as the five worked profiles, and
the same architectural evidence base discharges it. As the module states, this
is interpretive engineering guidance, not legal advice: confirm the mapping
with the institution's data-protection officer, information-security officer
and legal counsel before relying on it.

\endinput
